\section{Messages meet: the universe as one objective}
\label{sec:grmeet}

One edge makes a prediction.  A universe makes many at once: a dark
node may hear its calendar neighbour on each side and its index
besides, and the informers need not agree.  This section derives the
rule by which messages combine, first locally at one receiver, then
globally as a single least-squares objective for the whole universe
--- which will turn out to be \cref{ch:infer}'s update with many
unknowns at once.

\subsection{Two informers vote}

Each contract \cref{eq:gredge} is a quadratic penalty: believing the
relation at precision $p_{ij}$ means charging
$p_{ij}(\Theta_i-\beta_{ij}\Theta_j)^{2}$ for deviating from it, the
same precision-as-weight convention as every least-squares row since
Chapter~2.  Fix the informers for a moment and collect the terms that
involve one receiver $\Theta_i$:
\[
\sum_j p_{ij}\bigl(\Theta_i-\beta_{ij}\Theta_j\bigr)^{2}.
\]
This is a quadratic in $\Theta_i$, so complete the square.  Expanding
and gathering,
\[
\sum_j p_{ij}\Theta_i^{2}
-2\,\Theta_i\sum_j p_{ij}\beta_{ij}\Theta_j+\cdots
=\Bigl(\textstyle\sum_j p_{ij}\Bigr)
\Bigl(\Theta_i-\frac{\sum_j p_{ij}\,\beta_{ij}\Theta_j}
{\sum_j p_{ij}}\Bigr)^{2}+\cdots,
\]
where the dots are terms free of $\Theta_i$.  Read the completed
square:
\begin{equation}
\Theta_i \,\big|\, \text{informers}
\;\sim\;
\text{mean } \frac{\sum_j p_{ij}\,\beta_{ij}\Theta_j}{\sum_j p_{ij}},
\qquad
\text{variance } \frac{1}{\sum_j p_{ij}}.
\label{eq:grcond}
\end{equation}
Incoming messages --- each informer's move, mapped into the receiver's
units by its amplitude --- are precision-weighted and \emph{averaged};
the precisions themselves \emph{add}.  This is exactly the combination
rule of \cref{sec:infcombine} (two estimates of one number, weighted
by precision), meeting us again at a graph node.  Disagreement is
therefore a vote, never a sum: two opposing messages do not cancel
into silence --- they cancel into a confident zero.

\Cref{fig:grmeet}(a) runs three votes through the solver on the clean
ladder, with the 6M node dark and its neighbours lit at $-1$ (3M) and
$+1$ (1Y) volatility points.  First bar: amplitudes set to one and
equal trust --- the receiver lands at exactly
$\MacGrVoteCancel$, and its error bar is \emph{tighter} than a single
informer would leave it, because precisions added ($2p$) even as the
means cancelled.  Second bar: the short leg's trust raised to $3p$ ---
the vote goes to the louder voice,
$(3p\cdot(-1)+p\cdot(+1))/(4p)=\MacGrVoteOut$.  Third bar: equal
trust again but the \emph{calendar} amplitudes of
\cref{eq:grcalbeta} --- the same raw signals now read, in 6M units, as
$0.5\times(-1)=-0.5$ and $2\times(+1)=+2$, and their average is
$\MacGrVoteBeta$.  The receiver hears predictions in its own units,
not raw levels; equal-and-opposite cancellation was a statement about
amplitude-one edges, not a law.

\subsection{A chain: the mean crosses, the band pays}

Now let a message travel.  Chain two contracts,
$\Theta_B=\beta_1\Theta_A+\varepsilon_1$ and
$\Theta_C=\beta_2\Theta_B+\varepsilon_2$, with the noises independent
of each other and of $\Theta_A$.  Substitute the first into the
second:
\[
\Theta_C=\beta_2\beta_1\,\Theta_A
+\beta_2\,\varepsilon_1+\varepsilon_2 .
\]
Take the mean and the variance of both sides:
\begin{equation}
\E[\Theta_C]=\beta_2\beta_1\,\E[\Theta_A],
\qquad
\operatorname{Var}(\Theta_C)
=(\beta_2\beta_1)^{2}\operatorname{Var}(\Theta_A)
+\frac{\beta_2^{2}}{p_1}+\frac{1}{p_2}.
\label{eq:grchain}
\end{equation}
Amplitudes \emph{multiply} along the chain; relation variances
\emph{accumulate}, each mapped into the terminal node's units by the
amplitudes it still has to cross.  (In Chapter~10's electrical
picture: the hops are resistors in series --- each adds its
resistance, the variance, to the total --- while a receiver's several
informers in \cref{eq:grcond} are resistors in parallel, their
conductances, the precisions, adding.)  Nothing shrinks the mean:
distance taxes the band, never the mean.  A desk that wants remote
messages damped must therefore say so explicitly --- that is
\cref{sec:granchor}'s subject --- because the contracts alone will
never do it silently.

Panel~(b) of \cref{fig:grmeet} is \cref{eq:grchain} on a six-hop
amplitude-one chain with the far end lit at $+1$: the posterior mean
arrives at hop six at exactly $\MacGrHopMean$ --- undamped --- while
the band grows hop by hop, and the dotted curve, the closed-form
accumulation $\sqrt{\smash[b]{\mathcal{V}_{\mathrm{obs}}+k/p}}$ after
$k$ hops, lies on the solved band to machine precision (largest gap
\MacGrHopGap).  Compare this against \cref{fig:grcontract}: the two
figures are the same statement --- trust moves bands, amplitude moves
means --- read along a chain instead of across a dial.

\begin{figure}[htbp]
\centering
\paperfig{\textwidth}{fig_gr_meet.pdf}
\caption{Messages meeting.  (a)~The 6M receiver hears its lit
neighbours at $\mp1$ volatility point, three ways: equal trust at
amplitude one (exact cancellation, under a band tighter than either
informer alone would leave --- precisions added even as the means
cancelled); the short leg at $3p$ (the louder voice wins,
$\MacGrVoteOut$); equal trust under calendar amplitudes (the same
signals map to $-0.5$ and $+2$ in receiver units before averaging:
$\MacGrVoteBeta$).  Whiskers are $\pm1$ posterior sd.  (b)~A six-hop
chain, far end lit at $+1$: the mean crosses every hop undamped while
the band accumulates exactly the chain law \cref{eq:grchain} (dotted).}
\label{fig:grmeet}
\end{figure}

\subsection{The whole universe in one objective}

Local rules in hand, the global solve is now one display.  Collect
every belief the morning holds: the contracts, one quadratic each; the
lit observations, one row each; and (from \cref{sec:granchor}, zero
until then) an anchor per node.  The graph posterior is the minimizer
of
\begin{equation}
\boxed{\;
J(\Theta)
=\underbrace{\sum_{\text{edges}}
p_{ij}\bigl(\Theta_i-\beta_{ij}\Theta_j\bigr)^{2}}_{\text{contracts}}
\;+\;
\underbrace{\sum_{i}\kappa_i\,\Theta_i^{2}}_{\text{anchors}}
\;+\;
\underbrace{\sum_{s\ \text{lit}}
\frac{\bigl(\Theta_s-\hat\Theta_{\mathrm{obs},s}\bigr)^{2}}
{\mathcal{V}_{\mathrm{obs},s}}}_{\text{observations}} .
\;}
\label{eq:grjoint}
\end{equation}
Here $\hat\Theta_{\mathrm{obs},s}$ is lit node $s$'s measured
innovation --- panel~(b) of \cref{fig:gruniverse} --- and
$\mathcal{V}_{\mathrm{obs},s}$ its variance.  One point of accounting
before using it: an innovation is a \emph{difference} of two
estimates, the calibrated handle minus the transported baseline, and
the variance of a difference of independent estimates is the sum of
their variances (the same propagation Chapter~6 used for its error
bars): $\mathcal{V}_{\mathrm{obs},s}
=\mathcal{V}_{\mathrm{fit},s}+\mathcal{V}_{\mathrm{base},s}$.  A
node's baseline uncertainty enters the objective exactly once.  For a
lit node it is folded into $\mathcal{V}_{\mathrm{obs}}$, as just
shown.  For a dark node the mechanism is at the other end of the
pipeline: the published handle is baseline plus posterior innovation,
another difference-shaped sum of two estimates, so the published band
carries both variances,
$\mathcal{V}^{+}_i+\mathcal{V}_{\mathrm{base},i}$ --- the addition
made once, at publication.  Never both routes, and never twice.

Why is minimizing \cref{eq:grjoint} the right thing to do?  Because it
is \cref{ch:infer}'s update at universe scale.
Proposition~\ref{prop:infmap} proved the scalar case: minimizing a sum
of precision-weighted squares in one unknown \emph{is} the filter
update, and the objective's second derivative there came out as twice
the posterior precision.  Equation \cref{eq:grjoint} is the same
statement with a vector of unknowns: its matrix of second derivatives
is twice the universe's \emph{information matrix} $\mathcal{I}^{+}$
--- the same factor of two, now entrywise.  This is Chapter~10's
letter $\mathcal{I}$ at its third scale: there it priced one
functional's identifiability, then served as a whole fit's information
matrix; here it is a universe's.  The standard Gaussian facts
follow \citep{RasmussenWilliams2006,RueHeld2005}: the posterior mean
$\hat\Theta^{+}$ solves the normal equations
$\mathcal{I}^{+}\hat\Theta^{+}=b$, where $b$ collects the observation
targets (the anchors aim at zero, so they shape $\mathcal{I}^{+}$
only), and the \emph{marginal} posterior
variance of node $i$ is the $i$-th diagonal entry of the inverse,
$\mathcal{V}^{+}_i=(\mathcal{I}^{+})^{-1}_{ii}$.

\subsection{A two-node universe, by hand}

The smallest example already contains the two facts a reader must
carry out of this section.  One lit node $A$ (observed at
$\hat\Theta_{\mathrm{obs}}=d$ with variance $v$), one dark node $B$,
one contract $\Theta_B=\beta\Theta_A+\text{noise}$ at precision $p$.
The objective \cref{eq:grjoint} is
$(\Theta_A-d)^{2}/v+p(\Theta_B-\beta\Theta_A)^{2}$, and its
information matrix, differentiating twice with respect to
$(\Theta_A,\Theta_B)$, is
\[
\mathcal{I}^{+}
=\begin{pmatrix}1/v+p\beta^{2} & -p\beta\\[2pt] -p\beta & p\end{pmatrix},
\qquad
\det\mathcal{I}^{+}
=\frac{p}{v}+p^{2}\beta^{2}-p^{2}\beta^{2}
=\frac{p}{v},
\]
so, inverting the $2\times2$ matrix,
\[
(\mathcal{I}^{+})^{-1}
=\frac{v}{p}\begin{pmatrix}p & p\beta\\[2pt]
p\beta & 1/v+p\beta^{2}\end{pmatrix}
=\begin{pmatrix}v & \beta v\\[2pt]
\beta v & \dfrac{1}{p}+\beta^{2}v\end{pmatrix}.
\]
Read the two diagonal entries.  The lit node keeps its observation
variance $v$.  The dark node's marginal variance is
$1/p+\beta^{2}v$ --- the relation's own noise, plus the informer's
uncertainty mapped through the amplitude: precisely the chain law
\cref{eq:grchain} recovered from the joint solve.  (The mean,
solving the normal equations, is $\hat\Theta^{+}_A=d$,
$\hat\Theta^{+}_B=\beta d$: the contract at full amplitude.)

Now the caution.  The \emph{diagonal of the information matrix} at
$B$ reads $p$ --- call its reciprocal the \emph{conditional} variance,
because it answers ``how pinned is $B$ \emph{with $A$ held fixed}?''
The marginal variance $1/p+\beta^{2}v$ is larger, and it is the
correct one: $A$ is not fixed --- it is itself uncertain, and $B$
inherits that.  The gap between the two is exactly $\beta^{2}v$, the
informer's own doubt.  Reported confidence in this book is
\emph{always the marginal}; quoting the conditional overstates
certainty most at exactly the dark nodes whose informers are
themselves uncertain --- the nodes the graph exists for.

The objective is written and its output understood.  What remains is
to check the accounting --- that no configuration
mistake, no silent neighbour, no duplicated route can manufacture
information.  That is the next section.
